% Pollis drill-down: Localisation, Pollis in other languages, a draft (26 November 2026): what there is to translate,
% what VS Code already does, the language packs there are, the design of a Pollis language pack, of panels and wikis in
% each language and of toolboxes, how the reader chooses the language, the order of the work, and the open questions.
% Nothing in it is in Pollis yet. For whoever maintains Pollis, and for readers who would review a language.
% Compile with: latexmk -pdf workflow_localisation.tex
% Shared look for the one-page Pollis workflow handouts (same fonts and colours as the brochure).
\documentclass[10pt,a4paper]{article}

\usepackage[a4paper,left=17mm,right=17mm,top=14mm,bottom=16mm,footskip=8mm]{geometry}
\usepackage[T1]{fontenc}
\usepackage[default,semibold]{sourcesanspro}
\usepackage{microtype}
\usepackage{graphicx}
\usepackage[table]{xcolor}
\usepackage{tikz}
\usepackage[most]{tcolorbox}
\usepackage{fontawesome5}
\usepackage{enumitem}
\usepackage{fancyhdr}
\usepackage[hidelinks]{hyperref}

\graphicspath{{images/}}

% ---------------------------------------------------------------- colours (as in the brochure)
\definecolor{navy}{HTML}{1C3A5E}
\definecolor{pblue}{HTML}{3B78B5}
\definecolor{amber}{HTML}{F2B33D}
\definecolor{deepamber}{HTML}{D9902A}
\definecolor{pale}{HTML}{EEF3F9}
\definecolor{ink}{HTML}{2A2F36}
\definecolor{logonavy}{HTML}{173D63}
\definecolor{logodark}{HTML}{C98A1E}
\definecolor{logoamber}{HTML}{F2B446}
\definecolor{logocream}{HTML}{FEEFD2}
\definecolor{soft}{HTML}{5E6773}
\color{ink}

\setlength{\parindent}{0pt}
\setlength{\parskip}{3pt}
\setlist{nosep,leftmargin=*,itemsep=2pt}

\pagestyle{fancy}
\fancyhf{}
\renewcommand{\headrulewidth}{0pt}
\fancyfoot[L]{\footnotesize\color{soft}\textbf{Pollis}\enspace Bright Insights through Exploration, Modelling, Simulation and Optimisation}
\fancyfoot[R]{\footnotesize\color{soft}github.com/saragga/pollis/releases/latest}

% The Pollis logo, redrawn so it prints sharply at any size
\newcommand{\pollislogo}[1]{%
  \begin{tikzpicture}[x=#1/512,y=-#1/512,baseline=(current bounding box.center)]
    \fill[logonavy] (0,0) rectangle (512,512);
    \fill[logodark] (36,422) -- (104,256) -- (172,422) -- cycle;
    \fill[logoamber] (104,422) -- (203,180) -- (302,422) -- cycle;
    \fill[logocream] (203,422) -- (339,89) -- (476,422) -- cycle;
  \end{tikzpicture}}

% Top of the page: logo, name, kind of document (Workflow unless given, e.g. [Drill-Down]) and its subject;
% then the question and a one-line promise
\newcommand{\workflowhead}[4][Workflow]{%
  \thispagestyle{fancy}%
  \begin{minipage}[c]{13mm}\pollislogo{11mm}\end{minipage}%
  \begin{minipage}[c]{\dimexpr\linewidth-13mm\relax}
    {\color{navy}\Large\bfseries Pollis}\enspace{\color{soft}#1}\hfill
    {\color{deepamber}\small\bfseries\MakeUppercase{#2}}
  \end{minipage}\par\vspace{7mm}
  {\color{navy}\fontsize{22}{26}\selectfont\bfseries #3}\par\vspace{4pt}
  {\color{soft}\large #4}\par\vspace{5mm}}

% A menu path: \menu{Explore}{Access Data Libraries}... is written as \path{Explore \gt Access Data Libraries}
\newcommand{\gt}{\,{\color{deepamber}\bfseries ›}\,}
\newcommand{\menupath}[1]{{\footnotesize\bfseries\color{navy}#1}}

% A screenshot if images/<file> exists, otherwise an orange placeholder that says what to capture
\newlength{\shotheight}\setlength{\shotheight}{43mm}
\newcommand{\stepshot}[2]{%
  \IfFileExists{images/#1}{%
    \tcbox[enhanced, colback=black, colframe=black!70, boxrule=0.4pt, arc=2pt,
      left=0pt,right=0pt,top=0pt,bottom=0pt, boxsep=0pt, halign=center]{%
      \includegraphics[width=\dimexpr\linewidth-1pt\relax,height=\shotheight,keepaspectratio]{#1}}}{%
    \begin{tcolorbox}[enhanced, colback=white, colframe=orange!80, boxrule=0.6pt, arc=2pt,
      height=\shotheight, valign=center, halign=center, left=6pt, right=6pt]
      {\color{orange}\Large\faCamera}\par\smallskip
      {\color{orange}\footnotesize\bfseries Screenshot: #2}\par\smallskip
      {\color{orange!80!black}\scriptsize images/#1}
    \end{tcolorbox}}}

% One step: number, title, menu path, image file, what to capture, what it means
\newcommand{\step}[6]{%
  \begin{tcolorbox}
    {\color{deepamber}\Large\bfseries #1}\enspace{\color{navy}\large\bfseries #2}\par\vspace{1pt}
    \menupath{#3}\par\vspace{4pt}
    \stepshot{#4}{#5}\par\vspace{3pt}
    {\small #6}
  \end{tcolorbox}}

\newenvironment{steps}{%
  \begin{tcbraster}[raster columns=2, raster equal height=rows, raster column skip=4mm,
    raster row skip=4mm, enhanced, colback=pale, colframe=pale, boxrule=0pt, sharp corners,
    left=6pt,right=6pt,top=6pt,bottom=6pt, fontupper=\raggedright]}{\end{tcbraster}}

% Boxes for the bottom of the page
\newtcolorbox{pbox}[2][]{enhanced, colback=pale, colframe=amber, boxrule=0pt,
  leftrule=3pt, sharp corners, left=7pt, right=7pt, top=6pt, bottom=6pt,
  fontupper=\small\raggedright, before upper={{\color{navy}\normalsize\bfseries #2}\par\smallskip}, #1}
\newtcolorbox{nbox}[1][]{enhanced, colback=navy, colframe=navy, boxrule=0pt,
  sharp corners, left=8pt, right=8pt, top=6pt, bottom=6pt, coltext=white,
  fontupper=\small\raggedright, #1}

% A Next Steps item inside the bottom boxes
\newcommand{\nsitem}[2]{\item[\color{deepamber}\faChevronRight] \textbf{#1}: #2}

% Citations: \dcite{key}{Breiman (2001)} links to its entry \dref{key}{...} in the References, keyed as in pollis.bib
\newcommand{\dcite}[2]{\hyperlink{ref:#1}{\color{pblue}#2}}
\newcommand{\dref}[2]{\raisebox{1.1em}[0pt][0pt]{\hypertarget{ref:#1}{}}#2\par}
\newcommand{\ddoi}[1]{\href{https://doi.org/#1}{\color{pblue}doi:#1}}
\newcommand{\references}{\section*{References}\addcontentsline{toc}{section}{References}\footnotesize\raggedright\setlength{\parskip}{3pt}}

% The page of each heading, for Ask Pollis
% Pollis: the page of each heading, for Ask Pollis. A document that inputs this file, at the end of its preamble, writes
% <jobname>.labels beside its PDF, a line page|level|heading for each heading, the page counted from the first, the
% cover, as 1. compile_documents.jl copies it to ~/.pollis/documents, and ask_index.jl cuts the document into passages
% under its headings with it. \asklabel{1}{Title} marks a topic or a page, \asklabel{2}{Title} a heading inside it; the
% sections and subsections of a document that has them are marked as they go into its contents. A package vignette
% inputs it before its \documentclass (compile_vignette.jl), so the contents are taken over when the document begins,
% after the packages have taken them over. Read a second time, as by a copy of a document for Edit the LaTeX Source,
% which inputs it as resources/ask_labels when Pollis has put it in already (document_source.jl), it does nothing.
\ifdefined\asklabel\endinput\fi
\makeatletter
\newwrite\pollis@labels
\immediate\openout\pollis@labels=\jobname.labels
\newcommand{\asklabel}[2]{\write\pollis@labels{\the\ReadonlyShipoutCounter|#1|\unexpanded{#2}}}
\def\pollis@toc{toc}
\newcommand{\pollis@contentsline}[3]{%
  \def\pollis@file{#1}%
  \ifx\pollis@file\pollis@toc\asklabel{#2}{#3}\fi
  \pollis@addcontentsline{#1}{#2}{#3}}
% The headings of the New Methods style of the vignettes, \nmsection and \nmsub, which go into no contents
\def\pollis@asksection#1{\pollis@nmsection{#1}\asklabel{1}{#1}}
\def\pollis@asksub#1{\pollis@nmsub{#1}\asklabel{2}{#1}}
\AddToHook{begindocument/end}{%
  \let\pollis@addcontentsline\addcontentsline\let\addcontentsline\pollis@contentsline
  \ifdefined\nmsection\let\pollis@nmsection\nmsection\let\nmsection\pollis@asksection\fi
  \ifdefined\nmsub\let\pollis@nmsub\nmsub\let\nmsub\pollis@asksub\fi}
\makeatother



\usepackage{titlesec}
\usepackage{tabularx}
\usepackage{booktabs}
\usepackage{needspace}

\fancyhf{}
\fancyfoot[L]{\footnotesize\color{soft}\textbf{Pollis}\enspace Drill-down: Localisation\enspace·\enspace Draft, 26 November 2026}
\fancyfoot[R]{\footnotesize\color{soft}\thepage}

\titleformat{\section}{\color{navy}\large\bfseries}{\color{deepamber}\thesection}{0.6em}{}
\titleformat{\subsection}{\color{navy}\normalsize\bfseries}{\color{deepamber}\thesubsection}{0.5em}{}
\titlespacing*{\section}{0pt}{12pt}{4pt}
\titlespacing*{\subsection}{0pt}{8pt}{2pt}
\setlength{\parskip}{4pt}

\usepackage{xurl}
\newcommand{\file}[1]{{\small\nolinkurl{#1}}}
\newcolumntype{L}{>{\raggedright\arraybackslash}X}
\newcommand{\ui}[1]{\textbf{#1}}
\newcommand{\code}[1]{\texttt{\small #1}}

\begin{document}

\workflowhead[Drill-Down]{Localisation}
  {Pollis in other languages}
  {What there is to translate, what VS Code already does, and how Pollis's own menus, panels, wikis and toolboxes
   would follow the reader's language.}

\begin{nbox}
{\color{amber}\bfseries Drill-down, a draft.}\enspace
This drill-down is a design, not a description: nothing in it is in Pollis yet, and its settings, files and menu
items are proposals. Pollis today is in English, except VS Code's own menus, which follow a language pack.

It describes what there is to translate (\S\ref{sec:what}), what VS Code already does (\S\ref{sec:vscode}), the
language packs there are (\S\ref{sec:packs}), the design (\S\ref{sec:design}), how the reader chooses the language
(\S\ref{sec:choose}), languages in other scripts (\S\ref{sec:scripts}), review (\S\ref{sec:review}), the order of the
work (\S\ref{sec:order}) and the open questions (\S\ref{sec:questions}).
\end{nbox}

\tableofcontents

\section{What it is for}\label{sec:why}
\begin{itemize}
  \item \textbf{Readers work in their own language.} A student in Lyon or São Paulo reads a panel's key points and
    its wiki faster, and more surely, in French or Portuguese than in English.
  \item \textbf{The menus are half translated already.} With Microsoft's French pack, VS Code's menus are in French,
    while Pollis's own menus, pages and panels stay in English beside them.
  \item \textbf{The code stays as it is.} The Julia code of a panel, its comments and what it prints are code, which
    the reader edits and runs, and they stay in English, as the names of Julia's functions are.
\end{itemize}

\section{What there is to translate}\label{sec:what}
About 2.3 million words of Pollis's own text, besides VS Code's interface, counted in the sources on 26 November 2026.

\begin{center}\small
\renewcommand{\arraystretch}{1.2}
\begin{tabularx}{\linewidth}{@{}p{0.3\linewidth}p{0.24\linewidth}L@{}}
\toprule
\textbf{Layer} & \textbf{Size} & \textbf{How it is written today}\\
\midrule
VS Code's interface: menus, Settings, the Command Palette & about 19,000 strings & through \code{localize}; translated by
  Microsoft's language packs\\
Pollis's own interface: its pages, the frames of its panels, its menus & about 6,600 \code{localize} calls in
  Pollis's own files & through \code{localize}, as VS Code's; not translated\\
The built-in panels & 241 TOML files, 520,000 words with their code & plain English in each file\\
The built-in wikis & 1,762 pages, 960,000 words & plain English Markdown\\
The menus of the toolboxes & 15 toolboxes & plain English in each \file{package.json}\\
The toolboxes' panels and wikis & 121 panels, 420,000 words with their code; 686 wikis, 400,000 words & plain
  English\\
The documents of \ui{Help > Documentation} & the Tour, the Guides and the drill-downs & LaTeX; left out of this
  design\\
\bottomrule
\end{tabularx}
\end{center}

\section{What VS Code already does}\label{sec:vscode}
\begin{itemize}
  \item \textbf{Language packs.} A language pack is an extension with the translations of VS Code's strings, by their
    keys. \ui{Configure Display Language}, in the Command Palette, lists the packs installed and the language of the
    system, and sets the one chosen.
  \item \textbf{A restart.} VS Code reads its strings once, when it starts, so a new display language asks to restart.
  \item \textbf{An offer at the first start.} When the system's language has a pack, VS Code offers to install it and
    restart in that language. Pollis looks for the pack on Open VSX, as for every extension.
  \item \textbf{Extensions.} An extension's menus and settings take their titles from \file{package.nls.json}, and
    from \file{package.nls.fr.json} beside it in French; its code takes its strings from
    \file{l10n/bundle.l10n.fr.json}.
  \item \textbf{A version.} A pack serves the VS Code version it was made for and the later ones. Pollis is built on
    VS Code 1.117, so the newest pack it takes is 1.110.0, and the strings VS Code added after 1.110 show in English.
\end{itemize}

\section{The language packs there are}\label{sec:packs}
Microsoft publishes 13 packs on Open VSX, each with a version 1.110.0 that Pollis takes: French, German, Italian,
Spanish, Portuguese (Brazil), Russian, Polish, Czech, Turkish, Chinese (Simplified and Traditional), Japanese and
Korean. Other authors publish packs for other languages, on Open VSX or on the VS Code Marketplace, which Pollis's
\ui{Extensions} view does not search; a pack from there installs from its \file{.vsix} file, and its translations can
seed Pollis's own pack.

\begin{center}\small
\renewcommand{\arraystretch}{1.2}
\begin{tabularx}{\linewidth}{@{}p{0.2\linewidth}p{0.3\linewidth}p{0.17\linewidth}L@{}}
\toprule
\textbf{Language} & \textbf{Pack} & \textbf{VS Code strings} & \textbf{Where}\\
\midrule
French (and the 12 others above) & Microsoft's, 1.110.0 & 15,559 & Open VSX\\
Dutch & Nederlands Taalpakket, December 2024 & 11,678 & VS Code Marketplace\\
Arabic & Arabic-language.vscode-ar, October 2026 & 25,001 & VS Code Marketplace\\
Hindi & gdhillon.vscode-loc-hindi & 18,514 & Open VSX and the Marketplace\\
European Portuguese, Greek & none & --- & ---\\
\bottomrule
\end{tabularx}
\end{center}
The strings are those of the pack's \code{vscode} part, counted on 26 November 2026. The packs were made for different
versions of VS Code, so a pack made for a later version also has strings that Pollis's VS Code does not use.

\section{The design}\label{sec:design}
The code is written once; each language then costs its translation run and its review. Pollis follows VS Code's display
language, so a reader chooses French in one place.

\subsection{A Pollis language pack for each language}
The build already collects every \code{localize} string with its key. A script would merge Microsoft's pack for the
language with Pollis's own strings, translated, into one pack, \emph{Pollis Language Pack: Français}, published on
Open VSX with each Pollis release. One pack for each language, because VS Code uses one pack for a display language.
It also fills the strings VS Code added between 1.110 and 1.117, and a string not yet translated shows in English.

\subsection{One translation script for every layer}
\begin{itemize}
  \item \textbf{A model, with a glossary.} The script sends the texts in batches to a model the maintainer chooses,
    with the glossary of the language: \emph{Pollis}, \emph{Time Arrow}, the names of packages, data sets and menus of
    other programs stay as they are, and each language has fixed words for panel, step, wiki and toolbox. The
    instructions sent with the texts are a file in the repository, which anyone can read.
  \item \textbf{Checks on every answer.} The placeholders \code{\{0\}}, the \code{\&\&} of a menu's mnemonic, links,
    code spans, formulas and the \code{\{\{input\}\}} of a panel's code the same as in English; numbers that go into
    code, such as an input's default \code{1000}, left as they are. An answer that fails a check is asked for again.
  \item \textbf{A memory.} Each translation is kept with the English it came from, so a release translates only what
    changed, and a reviewer's correction is never translated over.
\end{itemize}

\subsection{Panels in each language}
Beside each panel's TOML, an overlay \file{l10n/fr/<panel>.toml} would hold only its text fields (title, intro,
labels, tooltips, key points, Next Steps), by the same ids, each with the English it translates. Code fields are never
in an overlay.
\begin{itemize}
  \item \textbf{Built-in panels.} \code{compile-toml} turns each panel's TOML into its data when Pollis is built; it
    would turn the overlays into the data of each language beside it, and the panel would take the language Pollis
    runs in.
  \item \textbf{Toolbox panels.} Pollis reads them from the toolbox when it opens them, and would read the overlay of
    the language beside each.
  \item \textbf{A stale translation.} When a panel's English changes, its overlay no longer matches: the field shows in
    English, and \code{compile-toml} and the panel validator list it, until it is translated again.
\end{itemize}

\subsection{Wikis in each language}
A copy \file{fr/<page>.md} beside each wiki page; the wiki view opens it when there is one, else the English page.
Citations, reference lists, formulas and code blocks stay as in English, and a copy records the English it translates,
as an overlay does.

\subsection{Toolboxes}
The titles in a toolbox's \file{package.json} become \code{\%keys\%}, with \file{package.nls.json} in English and
\file{package.nls.fr.json} beside it, which VS Code already reads. Its panels and wikis get overlays and copies as the
built-in ones. The Extension Agent and the packaging of a toolbox would write them when a toolbox is built.

\subsection{Ask Pollis}
The index and search of Ask Pollis stay in English in this design: its sources are the English texts, which a question
in another language would be searched against. Indexing each language is a later step.

\section{Choosing the language}\label{sec:choose}
\subsection{The display language: once, with a restart}
The display language is a choice for the whole of Pollis. It is made once, at the first start or in
\ui{Configure Display Language}, and asks to restart, because VS Code reads its strings when it starts. It needs no
place in the status bar: VS Code keeps the status bar for what belongs to the editor in front, which changes from one
to the next, such as a file's encoding, its line endings and its language mode.

\subsection{The language of a page: in the status bar}
A page's language is that kind of thing. As a LaTeX file declares its own encoding and its own languages, and the
status bar shows the encoding of the file in front, a status-bar item would show the language of the Pollis page or
wiki in front, when that page has a translation: \ui{Français}, or \ui{English} when it is read in English.
\begin{itemize}
  \item \textbf{A click switches that page} between its translation and the English, at once and without a restart,
    because a page's text comes from its TOML and its wiki, not from the strings read at the start. Its buttons and
    frame stay in the display language.
  \item \textbf{Why.} The translations are a model's until a reader reviews them; a reader who doubts a sentence reads
    the English in a click, and a reviewer compares the two.
  \item \textbf{A setting} would choose the default for every page: the display language, or English
    (\file{pollis.localisation.pages}). The item hides, as any status-bar item does, from its context menu.
  \item \textbf{A page in English only} shows no item: there is nothing to switch.
\end{itemize}

\section{Languages in other scripts}\label{sec:scripts}
The scripts themselves are no harder: Pollis runs on Chromium, which draws them all with the system's fonts. A
language is harder when VS Code has no pack for it, when it is written right to left, or when its plurals have more
forms than English.
\begin{itemize}
  \item \textbf{The same as French}: Japanese, Korean, and Chinese, simplified and traditional, have Microsoft's
    packs. Cantonese in Hong Kong is mostly written in standard written Chinese, so the traditional pack serves.
  \item \textbf{Plurals.} A noun after a number takes one form in English after 1 and another after the rest; Russian has
    three forms, for 1, for 2 to 4 and for 5 to 20, and Arabic six. VS Code's
    \code{localize} has none, so the strings with a count would choose their form with the browser's
    \code{Intl.PluralRules}, a string for each form.
  \item \textbf{No pack}: Greek, European Portuguese and Hindi (whose pack is another author's). The script would also
    translate VS Code's 19,000 strings, which are then Pollis's to keep up with each VS Code release.
  \item \textbf{Right to left}: Arabic. VS Code lays its workbench out left to right in every language, so its menus,
    sidebars and editors would stay as they are, with Arabic words in them. Pollis's own pages could be mirrored, page
    by page, with \code{dir="rtl"} and CSS's logical properties, which say start and end rather than left and right.
\end{itemize}

\section{Review}\label{sec:review}
\begin{itemize}
  \item \textbf{The interface}: a native speaker reads a sample of the strings of each language before its pack is
    published.
  \item \textbf{The scholarly text}: panel intros, key points and wikis are read page by page. Until a page is read, it
    says that it is a machine translation, with its model and date, and its status-bar item opens the English.
  \item \textbf{Corrections} go into the memory of \S\ref{sec:design}, so the next release keeps them; a reader sends
    one through \ui{Help > Report Issue}.
\end{itemize}

\section{The order of the work}\label{sec:order}
\begin{enumerate}
  \item \textbf{French, as a pilot}: the language pack, the \ui{Explore} menu and the Historical Geography and
    Archaeology toolbox, with the status-bar item and the checks.
  \item \textbf{German, Italian, Spanish and Portuguese}, with Microsoft's packs; then Dutch, from Nederlands
    Taalpakket's translations and Pollis's own.
  \item \textbf{The rest of Pollis's menus and the other toolboxes}, in the five languages.
  \item \textbf{Japanese, Korean, Chinese and Russian}, with Microsoft's packs and the plural forms.
  \item \textbf{Greek and Hindi}, with VS Code's strings translated by the script.
  \item \textbf{Arabic}, with its pages mirrored.
\end{enumerate}

\section{Open questions}\label{sec:questions}
\begin{enumerate}
  \item \textbf{Portuguese.} VS Code's only Portuguese pack is Brazilian. Should Pollis's own strings be Brazilian too,
    so that the menus read as one language, or European Portuguese, with VS Code's menus left in Brazilian?
  \item \textbf{The pilot.} \ui{Explore} and one toolbox, or only the menus and the panels' labels, with no wikis?
  \item \textbf{Review.} A reader for each language before it ships, or the scholarly text left as a machine
    translation, marked as one, until a reader comes?
  \item \textbf{Publishing.} The language packs on Open VSX under the \code{pollis} namespace, one for each language,
    installed from the \ui{Extensions} view as VS Code's are?
  \item \textbf{Ask Pollis.} An index for each language, or English sources found from a question translated first?
\end{enumerate}

\section{Where things would be}
\begin{center}\small
\renewcommand{\arraystretch}{1.2}
\begin{tabularx}{\linewidth}{@{}p{0.45\linewidth}L@{}}
\toprule
\textbf{File} & \textbf{What it holds}\\
\midrule
\file{build/pollis/l10n/} & The translation script, its instructions, the glossary and the memory of each language\\
\file{l10n/<language>/<panel>.toml}, beside each panel & A panel's text fields in the language\\
\file{<language>/<page>.md}, beside each wiki page & A wiki page in the language\\
\file{package.nls.<language>.json}, in each toolbox & A toolbox's menus in the language\\
\file{src/vs/workbench/contrib/localization/} & VS Code's \ui{Configure Display Language}, as it is\\
\bottomrule
\end{tabularx}
\end{center}

\end{document}
